% ======================================================================
% 致谢与参考文献
% ======================================================================
\section*{致谢}

作者在发展和理解格点规范理论的过程中受益于与许多人的交流。我感谢的人包括
J.~Bjorken、R.~P.~Feynman、M.~E.~Fisher、D.~Gross、J.~Kogut、L.~Susskind、
T.-M.~Yan，以及奥尔赛（Orsay）一个研讨班的成员。

\vspace{4mm}

\begin{center}
\footnotesize
\begin{minipage}{0.93\textwidth}
\raggedright
${}^{*}$本工作部分由美国国家科学基金会（National Science Foundation）资助。

\begin{enumerate}[leftmargin=1.6em,itemsep=1pt,parsep=0pt,topsep=2pt]
\item 参见，例如，J.~M. Cornwall and R.~E. Norton, Phys. Rev. D \textbf{8}, 3338 (1973); R.~Jackiw and K.~Johnson, \emph{ibid.} \textbf{8}, 2386 (1973); J.~Kogut and L.~Susskind, \emph{ibid.} \textbf{9}, 3501 (1974); D.~Amati and M.~Testa, Phys. Lett. \textbf{48B}, 227 (1974); G.~'t~Hooft, Nucl. Phys. B (to be published); P.~Olesen, Phys. Lett. \textbf{50B}, 255 (1974); A.~Chodos, R.~L.~Jaffe, K.~Johnson, C.~B.~Thorn, and V.~F.~Weisskopf, Phys. Rev. D \textbf{9}, 3471 (1974).
\item J.~Schwinger, Phys. Rev. \textbf{125}, 397 (1962); \textbf{128}, 2425 (1962).
\item J.~H. Lowenstein and J.~A. Swieca, Ann. Phys. (N.Y.) \textbf{68}, 172 (1971).
\item A.~Casher, J.~Kogut, and L.~Susskind, Phys. Rev. Lett. \textbf{31}, 792 (1973); Phys. Rev. D \textbf{10}, 732 (1974).
\item A.~Suri, Ph.D. thesis, Cornell University, 1969 (unpublished).
\item K.~G. Wilson and J.~Kogut, Phys. Rep. \textbf{12C}, 75 (1974).
\item S.~Coleman and E.~Weinberg, Phys. Rev. D \textbf{7}, 1888 (1973).
\item Y.~Nambu, in \emph{Symmetries and Quark Models}, proceedings of the International Conference on Symmetries and Quark Models, Wayne State University, 1969, edited by Ramesh Chand (Gordon and Breach, New York, 1970); L.~Susskind, Nuovo Cimento \textbf{69A}, 457 (1970); G.~Konisi, Prog. Theor. Phys. \textbf{48}, 2008 (1972); P.~Goddard, J.~Goldstone, C.~Rebbi, and C.~B.~Thorn, Nucl. Phys. \textbf{B56}, 109 (1972); J.-L.~Gervais and B.~Sakita, Phys. Rev. Lett. \textbf{30}, 716 (1973); S.~Mandelstam, Nucl. Phys. \textbf{B64}, 205 (1973).
\item R.~P. Feynman, Phys. Rev. \textbf{80}, 440 (1950).
\item 参见，例如，E.~S. Abers and B.~W. Lee, Phys. Rep. \textbf{9C}, 1 (1973).
\item 参见，例如，F.~A. Berezin, \emph{The Method of Second Quantization} (Academic, New York, 1966), pp.~52ff.
\item K.~G. Wilson, Phys. Rev. D \textbf{6}, 419 (1972).
\item M.~E. Fisher, Rep. Prog. Phys. \textbf{30}, 615 (1967).
\item 关于相对论性场论与临界点的关系见文献 6 第 X 节及其中引用的文献。综述见 R.~Brout, \emph{Phase Transitions} (Benjamin, New York, 1965), p.~8.
\item D.~Jasnow and M.~Wortis, Phys. Rev. \textbf{176}, 739 (1968).
\item K.~G. Wilson, Carg\`ese (1973) Lecture Notes, in preparation.
\item R.~Balian, J.~M. Drouffe, and C.~Itzykson, Phys. Rev. D (to be published).
\item J.~Kogut and L.~Susskind, Phys. Rev. D (to be published).
\item K.~Wilson, Cornell Report No.~CLNS-271 (to be published in the proceedings of the conference on Yang-Mills Fields, Marseille, 1974).
\item H.~P. Nielsen and P.~Olesen, Nucl. Phys. \textbf{B61}, 45 (1973); L.~J.~Tassie, Phys. Lett. \textbf{46B}, 397 (1973).
\end{enumerate}
\end{minipage}
\end{center}
